\chapter{Proposed Solution: Algorithm and Application}
\label{ch:solution}

% ============================================================
\section{Introduction}
\label{sec:sol-intro}
% ============================================================

This chapter presents the proposed system: a personalized, context-aware
\emph{next-POI engine} that is decoded into tourist \emph{itineraries}. Rather than
training a single monolithic itinerary generator, the design follows a deliberately
\emph{decoupled} philosophy. We first train a strong sequential model for the well-defined
next-POI prediction task (Phase~1), and then obtain itineraries by \emph{decoding} that
frozen model under inference-time constraints (Phase~2). This separation is not merely an
implementation convenience: as the experiments will show, decoding a strongly-supervised
next-POI model (\textbf{Frozen Rollout}) is a surprisingly hard baseline to beat with a
purpose-trained itinerary model (\textbf{Trained Pointer}).

The system therefore addresses the three research questions of Chapter~\ref{ch:introduction}
through the following concrete components: (i)~a data-preparation stage that curates a usable
check-in corpus and derives the spatial/temporal gap features the model consumes; (ii)~a
context-aware sequential engine combining a graph encoder, gap-feature context, a user
embedding, and a recurrent backbone (\textbf{RQ1, RQ2}); and (iii)~an inference-time decoding
procedure that turns next-POI scores into ordered routes (\textbf{RQ3}). The methods that the
thesis \emph{designed but did not evaluate} (explicit weather/time-of-day context, a
language-model persona encoder, embedding-based semantic pruning, and penalty-augmented
constrained decoding) are deferred to Chapter~\ref{ch:conclusion} as future work, so that this
chapter describes only what was actually built and measured.

The remainder of the chapter is organized as follows.
Section~\ref{sec:sol-overview} gives a high-level overview.
Section~\ref{sec:sol-data} describes data preparation.
Section~\ref{sec:sol-model} details the next-POI engine.
Section~\ref{sec:sol-decoding} presents Frozen Rollout and the Trained Pointer.
Section~\ref{sec:sol-experiments} specifies the experimental protocol, and
Section~\ref{sec:sol-results} reports and discusses the results.

% ============================================================
\section{System Overview}
\label{sec:sol-overview}
% ============================================================

The system operates in two phases, illustrated in Figure~\ref{fig:pipeline-overview}.

% Architecture of the proposed system (real, implemented pipeline).
%   Phase 1 — context-aware next-POI engine (GCN + Delta d/Delta t + user + GRU + head)
%   Phase 2 — itinerary decoding (Frozen Rollout / Trained Pointer) + benchmarks
% Requires: \usepackage{tikz} with libraries
%   positioning, arrows.meta, fit, shapes.geometric, calc
\begin{figure}[ht]
\centering
\resizebox{\textwidth}{!}{%
\begin{tikzpicture}[
    >=Stealth,
    font=\small,
    io/.style   ={draw, rounded corners=2pt, fill=gray!12, align=center, minimum height=0.85cm, inner sep=4pt},
    enc/.style  ={draw, rounded corners=2pt, fill=blue!12, thick, align=center, minimum height=0.85cm, inner sep=4pt},
    core/.style ={draw, rounded corners=2pt, fill=orange!15, thick, align=center, minimum height=0.85cm, inner sep=4pt},
    term/.style ={draw, rounded corners=2pt, fill=green!14, align=center, minimum height=0.85cm, inner sep=4pt},
    dec/.style  ={draw, rounded corners=2pt, fill=blue!8, align=center, minimum height=0.85cm, inner sep=4pt},
    merge/.style ={draw, circle, fill=white, inner sep=1pt, minimum size=0.5cm, font=\footnotesize},
    arr/.style  ={->, thick, gray!55!black},
]

% ---- Inputs / preprocessing ----
\node[io] (data) at (0,0) {Foursquare NYC\\check-ins};
\node[io, below=0.55cm of data] (prep) {Preprocessing:\\[-1pt]\footnotesize iter.\ $\ge\!10$ filter,\\[-1pt]\footnotesize 24\,h sessions, $\Delta d,\Delta t$};

% ---- Three encoder branches ----
\node[io,  right=1.5cm of data.east, yshift=1.35cm] (graph) {Hybrid POI graph\\[-1pt]\footnotesize co-visit $\cup$ 10-NN};
\node[io,  below=0.5cm of graph] (ctxin) {$\Delta d,\ \Delta t$\\[-1pt]\footnotesize per step};
\node[io,  below=0.5cm of ctxin] (userid) {User ID};

\node[enc, right=1.2cm of graph]  (gcn)  {2-layer\\GCN};
\node[enc, right=1.2cm of ctxin]  (cmlp) {Context\\MLPs};
\node[enc, right=1.2cm of userid] (uemb) {User\\embedding $e_u$};

% ---- Merge POI features + context, feed GRU ----
\node[merge, right=0.7cm of gcn, yshift=-0.7cm] (cat1) {$\oplus$};
\node[core, right=0.8cm of cat1] (gru) {GRU};
\node[merge, right=0.7cm of gru, yshift=-1.4cm] (cat2) {$\oplus$};
\node[core, right=0.8cm of cat2] (head) {MLP head\\+ softmax};
\node[term, right=0.9cm of head] (nextpoi) {Next-POI\\ranking over $|\mathcal{V}|$\\[-1pt]\footnotesize HR@k, NDCG, MRR};

% ---- Arrows: Phase 1 ----
\draw[arr] (data) -- (prep);
\draw[arr] (prep.east) -- (graph.west);
\draw[arr] (prep.east) -- (ctxin.west);
\draw[arr] (prep.east) -- (userid.west);
\draw[arr] (graph) -- (gcn);
\draw[arr] (ctxin) -- (cmlp);
\draw[arr] (userid) -- (uemb);
\draw[arr] (gcn.east)  -- (cat1);
\draw[arr] (cmlp.east) -- (cat1);
\draw[arr] (cat1) -- (gru);
\draw[arr] (gru.east) -- (cat2);
\draw[arr] (uemb.east) -| (cat2);
\draw[arr] (cat2) -- (head);
\draw[arr] (head) -- (nextpoi);

% ---- Phase 2: itinerary decoding ----
\node[dec, below=4.7cm of gru, xshift=-2.2cm] (rollout)
  {\textbf{Frozen Rollout}\\[-1pt]\footnotesize frozen engine + no-revisit mask,\\[-1pt]\footnotesize reserved end, greedy / beam};
\node[dec, left=1.0cm of rollout] (pointer)
  {Trained Pointer\\[-1pt]\footnotesize GCN+pointer, end-to-end};
\node[term, right=1.0cm of rollout] (itin)
  {Ordered itinerary\\[-1pt]\footnotesize pairs-F1, set-F1};
\node[dec, right=0.8cm of itin] (flickr)
  {Flickr Benchmark\\[-1pt]\footnotesize Chen-2016 protocol};

\draw[arr, blue!55!black, densely dashed] (head.south) -| (rollout.north);
\draw[arr] (pointer) -- (rollout);
\draw[arr] (rollout) -- (itin);
\draw[arr] (itin) -- (flickr);

% ---- Phase containers (dashed fit boxes, no fill so they sit over the main layer) ----
\node[draw, dashed, rounded corners, inner sep=11pt, label={[anchor=north west,font=\footnotesize\bfseries,text=blue!50!black,xshift=2pt,yshift=-2pt]north west:Phase 1 --- context-aware next-POI engine},
      fit=(data)(graph)(gcn)(gru)(head)(nextpoi)(userid)] (p1box) {};
\node[draw, dashed, rounded corners, inner sep=9pt, label={[anchor=north west,font=\footnotesize\bfseries,text=orange!55!black,xshift=2pt,yshift=-2pt]north west:Phase 2 --- itinerary decoding \& benchmarks},
      fit=(pointer)(rollout)(itin)(flickr)] (p2box) {};

\end{tikzpicture}%
}
\caption{Architecture of the proposed system. \textbf{Phase~1} (the next-POI engine): a
two-layer GCN over a hybrid POI graph (training-set co-visitation edges $\cup$ geographic
10-nearest-neighbour edges) produces POI features; per-step spatial/temporal gaps
$(\Delta d,\Delta t)$ are lifted by small MLPs; a learnable user embedding $e_u$ encodes
identity. The concatenated POI-feature/context sequence is read by a GRU, whose final state
is concatenated with $e_u$ and scored by an MLP head over the full POI vocabulary
($\mathrm{HR@}k$, $\mathrm{NDCG@}k$, $\mathrm{MRR}$). \textbf{Phase~2} (itinerary
recommendation): the \emph{frozen} engine is decoded into ordered routes by \emph{Frozen
Rollout} (no-revisit mask, reserved end POI, greedy or beam), compared against an end-to-end
\emph{Trained Pointer}; both are scored with set-F1 / pairs-F1 and validated on the Flickr
datasets under the Chen-2016 protocol.}
\label{fig:pipeline-overview}
\end{figure}


\textbf{Phase~1 --- the next-POI engine.} Raw Foursquare check-ins are filtered and grouped
into per-user sessions, and each consecutive visit pair is annotated with its spatial gap
$\Delta d$ (haversine distance) and temporal gap $\Delta t$ (elapsed time). A two-layer Graph
Convolutional Network (GCN) over a hybrid POI graph produces a feature vector for every POI;
the GCN features and the per-step context $(\Delta d,\Delta t)$ form the input to a GRU, whose
final state is combined with a learnable user embedding and scored by a multilayer-perceptron
(MLP) head over the full POI vocabulary. The engine is trained for single-step next-POI
prediction and evaluated with ranking metrics ($\mathrm{HR@}k$, $\mathrm{NDCG@}k$,
$\mathrm{MRR}$).

\textbf{Phase~2 --- itinerary decoding.} Given a query (a user, a start POI, an optional end
POI, and a desired length), \emph{Frozen Rollout} rolls the frozen engine forward, masking
already-visited POIs and reserving the final step for the end POI, to produce an ordered route.
As an integrated alternative, a \emph{Trained Pointer} replaces the MLP head with an
attention-style pointer and is trained end-to-end on whole trajectories. Itineraries are scored
with set-F1 and order-aware pairs-F1, and---because the Foursquare numbers are not on the same
scale as the published trip-recommendation literature---are re-validated on the Flickr
photo-trajectory datasets under the exact protocol of Chen et al.~\cite{chen2016learning}.

% ============================================================
\section{Data Preparation}
\label{sec:sol-data}
% ============================================================

\subsection{Filtering and sessionization}
\label{sec:sol-data-filter}

We use the standard Foursquare NYC check-in dataset of Yang et al.~\cite{yang2015foursquare}
($\sim$227K check-ins, $\sim$1.08K users, $\sim$38K venues). As is standard in next-POI
research, the raw data is reduced to a denser, more learnable core by an \emph{iterative
$k$-core filter}: users with fewer than ten check-ins and POIs with fewer than ten check-ins
are removed, and the procedure is repeated until the set is stable (in practice a few
iterations). This retains roughly $5{,}100$ POIs and $\sim$1.0K users. We deliberately keep
this filter simple and reproducible; a tourism-specific semantic curation of the catalogue is
discussed as future work (Chapter~\ref{ch:conclusion}).

User and POI identifiers are re-indexed to contiguous integers. Check-ins are then grouped
into \textbf{sessions}: a gap of more than 24~hours between successive check-ins of the same
user starts a new session, and sessions of length one (which admit no prediction) are dropped.

\subsection{Spatial and temporal gap features}
\label{sec:sol-data-context}

For every consecutive pair of check-ins within a session we compute two context scalars: the
\textbf{spatial gap} $\Delta d_t$ (haversine distance in km from the previous POI, clipped to
$[0,100]$) and the \textbf{temporal gap} $\Delta t_t$ (elapsed time in hours, clipped to
$[0,24]$). At the first step of a session $\Delta d_1=\Delta t_1=0$. These two gaps are the
\emph{only} contextual signals the engine consumes; they capture how far and how long a user
moved between visits, which is the spatio-temporal regularity exploited by ST-RNN-style models
\cite{liu2016strnn}. Time-of-day and weather are not modelled.

\subsection{Hybrid POI graph}
\label{sec:sol-data-graph}

To let the model reason about POIs beyond their individual identity, we build a single
undirected \textbf{hybrid graph} $G=(\mathcal{V},\mathcal{E}_{\mathrm{cov}}\cup
\mathcal{E}_{\mathrm{geo}})$ on the curated catalogue:
\begin{itemize}[leftmargin=*]
  \item \textbf{Co-visitation edges} $\mathcal{E}_{\mathrm{cov}}$: an edge $(p_i,p_j)$ is added
        when at least three consecutive $p_i\!\rightarrow\!p_j$ transitions occur within the
        \emph{training} sessions, encoding behavioural proximity;
  \item \textbf{Geographic edges} $\mathcal{E}_{\mathrm{geo}}$: each POI is connected to its
        ten nearest neighbours by haversine distance (computed with a ball tree), encoding
        spatial proximity.
\end{itemize}
The two edge sets are unioned and symmetrized. Co-visitation edges are derived from training
data only, so no test information leaks into the graph.

\subsection{Splitting}
\label{sec:sol-data-split}

Sessions are split \textbf{chronologically per user} into 70\%/10\%/20\% train/validation/test,
so that no future session of a user is used to predict its past. This split is shared by the
next-POI evaluation and (its test portion) by the NYC itinerary evaluation.

% ============================================================
\section{The Next-POI Engine}
\label{sec:sol-model}
% ============================================================

The engine maps a session prefix $S_u=(p_1,\dots,p_T)$ and its user $u$ to a probability
distribution over the next POI. It has four learnable parts: a graph encoder, a context
encoder, a user embedding, and a recurrent scorer. Default dimensions are
$d_p\!=\!128$ (POI), $d_u\!=\!64$ (user), $d_c\!=\!32$ (context), and $d_h\!=\!128$ (GRU hidden).

\subsection{Graph encoder}
\label{sec:sol-model-gcn}

A two-layer GCN over the hybrid graph transforms a learnable POI embedding table
$E^{(0)}\in\mathbb{R}^{|\mathcal{V}|\times d_p}$ into graph-aware POI features:
\begin{equation}
  E^{(l+1)} = \sigma\!\left(\tilde{D}^{-1/2}\,\tilde{A}\,\tilde{D}^{-1/2}\,E^{(l)}\,W^{(l)}\right),
  \qquad l = 0,1,
  \label{eq:gcn}
\end{equation}
where $\tilde{A}=A+I$ is the adjacency with self-loops, $\tilde{D}$ its degree matrix,
$\sigma$ a ReLU non-linearity (with dropout between layers), and $W^{(l)}$ the layer weights.
The POI feature used downstream is $h_p^{\mathrm{GCN}}=E^{(2)}_p\in\mathbb{R}^{d_p}$. Because the
graph is fixed, the full forward pass of Eq.~\eqref{eq:gcn} is computed \emph{once} per batch
(and once per decode) and indexed for each POI in a sequence.

\subsection{Context encoder}
\label{sec:sol-model-context}

The two gap scalars are lifted to a $d_c$-dimensional context vector by two small MLPs and
concatenated:
\begin{equation}
  c_t = \big[\,\mathrm{MLP}_d(\Delta d_t)\;\Vert\;\mathrm{MLP}_t(\Delta t_t)\,\big]
        \in\mathbb{R}^{d_c}.
  \label{eq:context}
\end{equation}

\subsection{User embedding}
\label{sec:sol-model-user}

User identity is encoded by a learnable embedding lookup $e_u\in\mathbb{R}^{d_u}$. This is a
collaborative, identity-based representation learned from the user's own visit history during
training; it is not derived from a persona description. The contribution of this embedding is
isolated by a dedicated ablation in Section~\ref{sec:sol-results-pers}.

\subsection{Recurrent scorer and output}
\label{sec:sol-model-gru}

At each step the GCN feature and the context are concatenated,
$x_t=[\,h_{p_t}^{\mathrm{GCN}}\;\Vert\;c_t\,]\in\mathbb{R}^{d_p+d_c}$, and the sequence
$X=(x_1,\dots,x_T)$ is read by a GRU. Its final hidden state $h_T\in\mathbb{R}^{d_h}$ is
concatenated with the user embedding and scored by an MLP head over the whole vocabulary:
\begin{align}
  h_T &= \mathrm{GRU}(x_1,\dots,x_T), \qquad z = [\,h_T\;\Vert\;e_u\,], \label{eq:gru}\\
  \hat{y} &= W_2\,\mathrm{ReLU}(W_1 z + b_1) + b_2 \in \mathbb{R}^{|\mathcal{V}|}, \label{eq:head}\\
  P(p_{T+1}&=p_j\mid S_u,u) = \frac{\exp(\hat{y}_j)}{\sum_k \exp(\hat{y}_k)}.
\end{align}
The model is trained by minimizing the next-POI cross-entropy over all session prefixes:
\begin{equation}
  \mathcal{L} = -\frac{1}{N}\sum_{i=1}^{N}\log P\!\left(y^{(i)}\mid S_u^{(i)},u^{(i)}\right),
  \label{eq:loss}
\end{equation}
with the Adam optimizer, gradient clipping, dropout regularization, and early stopping on
validation $\mathrm{HR@}10$. A session of length $L$ contributes $L-1$ prefix/target training
pairs, giving dense supervision (on the order of $10^5$ examples on NYC).

% ============================================================
\section{From Next-POI Scores to Itineraries}
\label{sec:sol-decoding}
% ============================================================

\subsection{Problem and query}
\label{sec:sol-decoding-query}

Itinerary recommendation takes a query $q=(u,s,e,K)$---a user $u$, a start POI $s$, an
optional end POI $e$, and a length budget $K$---and returns an ordered, loop-free route
$\hat{Y}=(s,p_2,\dots,e)$ of length $K$. Following the Chen-2016 convention
\cite{chen2016learning}, the start, the end, and the length are \emph{given}, and quality is
how well $\hat{Y}$ recovers the user's actual visited sequence $Y$. We use the length budget
(read directly as the ground-truth session length) as the headline mode, because it needs no
invented travel-speed or dwell-time constants and is therefore fully reproducible.

\subsection{Frozen Rollout (the headline method)}
\label{sec:sol-decoding-rollout}

The key observation is that the Phase-1 engine already scores a distribution over the next POI
for any partial route. An itinerary is what we obtain by \emph{rolling that scorer out} under
two constraints it never saw at training time:
\begin{enumerate}[label=(\roman*),leftmargin=*]
  \item \textbf{No-revisit mask}: the logit of every already-visited POI is set to $-\infty$,
        enforcing loop-free routes \cite{menon2017revisits};
  \item \textbf{Reserved end}: at the final step the route is forced to the queried end POI
        $e$ (the discrete analogue of the orienteering start/end constraint).
\end{enumerate}
\emph{Nothing in the trained model changes}---the same checkpoint serves both tasks; the entire
itinerary capability is an inference-time decoder. Two decoders are used: \textbf{greedy}
(take the \texttt{argmax} of the masked logits at each step) as the deterministic floor, and
\textbf{beam($b$)} (keep the $b$ partial routes with the highest summed transition
log-probability). At decode time $\Delta d$ is the real haversine distance between consecutive
route POIs and $\Delta t$ is a single documented constant; in length-budget mode this constant
affects only the context features, not the stop rule. Crucially, this headline decoder uses
\emph{no} travel-cost penalty and \emph{no} diversity bonus---augmenting the beam objective
with such penalties is left to future work (Chapter~\ref{ch:conclusion}).

\subsection{Trained Pointer (the integrated alternative)}
\label{sec:sol-decoding-pointer}

To test whether a model trained \emph{directly} for itineraries beats the decoded engine, we
also implement a pointer-network decoder \cite{vinyals2015pointer}. It reuses the GCN encoder
and user embedding but replaces the fixed MLP head with an \textbf{inner-product pointer}:
the decoder state $h_t$ is projected and scored against every POI feature,
$\mathrm{logit}_v=\langle W_o[h_t\Vert e_u],\,h_v^{\mathrm{GCN}}\rangle$, tying the output
directly to the graph signal. It is trained with teacher-forced sequence cross-entropy over
\emph{whole} length-$\ge\!3$ trajectories (one example per session), with the no-revisit mask
applied only at decode time and early stopping on validation pairs-F1. A context-aware variant
additionally feeds the $(\Delta d,\Delta t)$ context at each decoder step. Because each session
yields a single whole-trajectory example, the pointer sees far fewer updates than the
prefix-rich next-POI engine---a supervision-density gap that the results make visible.

% ============================================================
\section{Experimental Setup}
\label{sec:sol-experiments}
% ============================================================

\subsection{Two tasks, two benchmarks}
\label{sec:sol-exp-tasks}

This thesis evaluates two \emph{formally distinct} sub-tasks, each on its own field-standard
benchmark, following the surveys of Lim et al.~\cite{lim2018survey} and
Halder et al.~\cite{halder2024survey}:
\begin{itemize}[leftmargin=*]
  \item \textbf{Next-POI prediction} on Foursquare NYC, with full-vocabulary
        $\mathrm{HR@}k$/$\mathrm{NDCG@}k$/$\mathrm{MRR}$ (the protocol of GETNext
        \cite{yang2022getnext} and LLM4POI \cite{li2024llm4poi});
  \item \textbf{Itinerary recommendation} with set-F1 and pairs-F1, reported both internally on
        Foursquare NYC and---for literature comparison---on the Flickr photo-trajectory cities
        under the Chen-2016 protocol \cite{chen2016learning}.
\end{itemize}
The two tasks' metrics are \textbf{not on a common scale and are never cross-compared}: a
next-POI $\mathrm{HR@}1$ and an itinerary pairs-F1 measure different things, and absolute scores
across datasets with different POI-vocabulary sizes are not comparable
\cite{krichene2020sampled}. This is why the NYC itinerary numbers (Section~\ref{sec:sol-results-nyc})
and the Flickr numbers (Section~\ref{sec:sol-results-flickr}) live in separate tables.

\subsection{Metrics}
\label{sec:sol-exp-metrics}

\textbf{Next-POI:} $\mathrm{HR@}k$ (the true next POI is ranked in the top $k$ of the full
vocabulary), $\mathrm{NDCG@}k$, and $\mathrm{MRR}$, all without sampled negatives.

\textbf{Itinerary:} for a route, let its ordered-pair set be
$P(T)=\{(T_i,T_j):i<j\}$. Then
$\text{pairs-F1}$ is the F1 over $P(\hat{Y})$ vs $P(Y)$ (a pair counts only if both POIs appear
in both routes \emph{in the same order}), and $\text{set-F1}$ is the order-agnostic F1 over the
visited \emph{sets}. We report both, plus the exact-match rate. The pairs-F1 implementation is
pinned by unit tests against the Chen-2016 reference, and---per the faithfulness rule below---a
\texttt{Random} baseline reproduces Chen's published \texttt{Random} within noise. Following the
literature, itinerary metrics are reported on \textbf{length-$\ge\!3$, loop-free} trajectories,
where pairs-F1 has dynamic range (length-2 queries are trivial $s\!\to\!e$ recoveries).

\subsection{Baselines}
\label{sec:sol-exp-baselines}

For \emph{next-POI}, we position the engine against the published tier from the LLM4POI
benchmark table (LSTM, STGCN \cite{yu2018stgcn}, STAN \cite{luo2021stan}, GETNext
\cite{yang2022getnext}, STHGCN \cite{yan2023sthgcn}, LLM4POI \cite{li2024llm4poi}). For the
\emph{itinerary} task, the NYC study contrasts Frozen Rollout against the Trained Pointer, and
the Flickr Benchmark adds classical baselines (\texttt{Random}, \texttt{PoiPopularity},
\texttt{Markov}, \texttt{MarkovPath}) and the published methods PoiRank/Rank+Markov
\cite{chen2016learning}, DeepTrip \cite{gao2019deeptrip}, CTLTR \cite{zhou2021ctltr}, SelfTrip
\cite{gao2022selftrip}, and AR-Trip \cite{shu2024artrip}.

% ============================================================
\section{Results and Discussion}
\label{sec:sol-results}
% ============================================================

\subsection{Next-POI accuracy (Phase~1)}
\label{sec:sol-results-nextpoi}

Table~\ref{tab:nextpoi} reports the engine's full-vocabulary ranking accuracy on Foursquare
NYC. The model reaches $\mathrm{HR@}1=0.187$, i.e.\ the exact next venue is the top prediction
out of $\sim$5{,}100 candidates almost one time in five, with $\mathrm{MRR}=0.316$.

\begin{table}[ht]
\centering
\caption{Next-POI prediction accuracy of the proposed engine on Foursquare NYC
(full vocabulary, test set).}
\label{tab:nextpoi}
\begin{tabular}{lcccccc}
\toprule
Model & HR@1 & HR@5 & HR@10 & NDCG@5 & NDCG@10 & MRR \\
\midrule
Ours (GCN+GRU+user) & 0.187 & 0.476 & 0.588 & 0.339 & 0.375 & 0.316 \\
\bottomrule
\end{tabular}
\end{table}

To position this honestly, Table~\ref{tab:nextpoi-tier} places our $\mathrm{HR@}1$ within the
published tier (numbers as collated in the LLM4POI benchmark). The engine sits in the
\emph{LSTM/STGCN band}: clearly learning (well above an LSTM and on par with a spatio-temporal
GCN) but below the attention/transformer SOTA (STAN, GETNext, STHGCN) and far below the
LLM-based LLM4POI. This is the expected, defensible position for a compact GCN+GRU model and is
stated as such rather than overclaimed.

\begin{table}[ht]
\centering
\caption{$\mathrm{HR@}1$ (equivalently $\mathrm{Acc@}1$) tier on Foursquare NYC, as collated in
the LLM4POI benchmark. Our engine is in bold.}
\label{tab:nextpoi-tier}
\begin{tabular}{lc}
\toprule
Model & HR@1 \\
\midrule
LSTM & 0.130 \\
STGCN~\cite{yu2018stgcn} & 0.180 \\
\textbf{Ours (GCN + GRU + user)} & \textbf{0.187} \\
STAN~\cite{luo2021stan} & 0.220 \\
GETNext~\cite{yang2022getnext} & 0.240 \\
STHGCN~\cite{yan2023sthgcn} & 0.270 \\
LLM4POI~\cite{li2024llm4poi} & 0.340 \\
\bottomrule
\end{tabular}
\end{table}

\subsection{Itinerary on NYC: Frozen Rollout vs.\ Trained Pointer}
\label{sec:sol-results-nyc}

Table~\ref{tab:nyc-itin} reports the headline integration experiment on the NYC test sessions
($n=2{,}880$ length-$\ge\!3$ trajectories). The result is deliberately counter-intuitive:
\textbf{decoding the frozen next-POI engine (decoupled) beats the purpose-trained pointer
(integrated)} by about $0.03$ pairs-F1 ($-10.5\%$ for the pointer). Beam search barely improves
on greedy ($+0.0015$ pairs-F1), which shows the limitation is the engine's myopia, not the
decoder: smarter search cannot recover information the scorer never had.

\begin{table}[ht]
\centering
\caption{Itinerary recommendation on Foursquare NYC (length-$\ge\!3$ test, $n=2{,}880$). Higher
is better; itinerary metrics are \emph{not} comparable to the next-POI metrics of
Table~\ref{tab:nextpoi}.}
\label{tab:nyc-itin}
\begin{tabular}{llccc}
\toprule
Method & Type & pairs-F1 & set-F1 & exact \\
\midrule
\textbf{Frozen Rollout} (greedy) & decoupled & \textbf{0.289} & 0.609 & 0.054 \\
Frozen Rollout (beam~3) & decoupled & 0.290 & 0.610 & 0.057 \\
Trained Pointer v1 (no context) & integrated & 0.259 & 0.578 & 0.043 \\
Trained Pointer v2 (\,+\,context) & integrated & 0.261 & --- & --- \\
\bottomrule
\end{tabular}
\end{table}

The pointer trained cleanly (validation pairs-F1 rising to a genuine peak before early
stopping), so this is a real negative result, not a bug. Three factors explain why the simpler
route wins: (i)~\textbf{supervision density}---the engine learns from $\sim$$10^5$ prefix/target
pairs (every prefix of every session), whereas the pointer sees only one whole-trajectory
example per session; (ii)~\textbf{leaner inputs}---the v1 pointer conditions only on the
previous POI's graph feature, which is why adding the $(\Delta d,\Delta t)$ context in v2
recovers a small amount; and (iii)~\textbf{converged vs.\ from-scratch}---Frozen Rollout reuses
a fully-trained engine while the pointer trains its encoder from random initialization. The
takeaway for the thesis is that \emph{naively training a dedicated itinerary model does not
automatically beat cleverly decoding a strong next-POI model}---a nuance to the integration
narrative of Halder et al.~\cite{halder2025drlir}.

A second, important caveat is that these NYC pairs-F1 values ($\approx0.26$--$0.29$) are
\textbf{not comparable} to the published trip-recommendation numbers ($\approx0.6$--$0.85$):
the published works evaluate on Flickr cities of $\sim$25--90 POIs, whereas NYC has $\sim$5{,}100
POIs, so the chance level differs by two orders of magnitude. The gap is a dataset/vocabulary
artifact, not a quality gap---which is precisely why the next subsection re-runs the task on the
literature's own benchmark.

\subsection{The Flickr Benchmark (literature-comparable)}
\label{sec:sol-results-flickr}

To place the itinerary task on the same scale as the published literature, we re-run it on the
five Flickr photo-trajectory cities under the \emph{exact} Chen-2016 protocol
\cite{chen2016learning}: leave-one-trajectory-out cross-validation, the query gives the first
and last POI plus the desired length, and evaluation is on length-$\ge\!3$ loop-free
trajectories. Table~\ref{tab:flickr-data} lists the datasets.

\begin{table}[ht]
\centering
\caption{Flickr photo-trajectory datasets (canonical Chen-2016 files). Evaluation is on the
length-$\ge\!3$ subset.}
\label{tab:flickr-data}
\begin{tabular}{lcccc}
\toprule
City & \#POIs & \#users & \#trajectories & \#traj.\ (len $\ge 3$) \\
\midrule
Toronto   & 29 & 1{,}395 & 6{,}057 & 335 \\
Osaka     & 27 &   450 & 1{,}115 &  47 \\
Glasgow   & 27 &   601 & 2{,}227 & 112 \\
Edinburgh & 28 & 1{,}454 & 5{,}028 & 634 \\
Melbourne & 88 & 1{,}000 & 5{,}106 & 442 \\
\bottomrule
\end{tabular}
\end{table}

\paragraph{Protocol faithfulness.} A \texttt{Random} baseline has no modelling choices, so it
must reproduce Chen's \texttt{Random} if and only if the data, protocol, and metric are
identical. Table~\ref{tab:flickr-faith} shows they match within noise (and our
\texttt{PoiPopularity} matches Chen's near-exactly on the data-rich cities), which validates
that our harness is protocol-faithful---a prerequisite for any comparison to the literature.

\begin{table}[ht]
\centering
\caption{Faithfulness check (pairs-F1): our reproductions vs.\ Chen 2016.}
\label{tab:flickr-faith}
\begin{tabular}{llccccc}
\toprule
Method & Source & Toronto & Osaka & Glasgow & Edinburgh & Melbourne \\
\midrule
\multirow{2}{*}{Random} & ours & 0.298 & 0.301 & 0.301 & 0.270 & 0.227 \\
                        & Chen 2016 & 0.310 & 0.304 & 0.320 & 0.261 & 0.248 \\
\addlinespace
\multirow{2}{*}{PoiPopularity} & ours & 0.443 & 0.413 & 0.510 & 0.439 & 0.320 \\
                               & Chen 2016 & 0.384 & 0.365 & 0.507 & 0.436 & 0.316 \\
\bottomrule
\end{tabular}
\end{table}

\paragraph{Our results.} Table~\ref{tab:flickr-ours} reports our classical baselines and the
learned pointer. The headline is that \textbf{our pairs-F1 now spans $0.23$--$0.59$, on the
published $0.26$--$0.85$ scale}, versus $0.26$--$0.29$ on NYC---confirming that the low NYC
numbers were a property of the data and protocol, not the method. Among our methods the
first-order \texttt{Markov} transition baseline is strongest, and the light from-scratch pointer
(0.31--0.49) reaches the scale and beats \texttt{Random} comfortably but trails \texttt{Markov}
on every city---the same supervision-density lesson seen on NYC.

\begin{table}[ht]
\centering
\caption{Our itinerary results on the Flickr cities (pairs-F1, leave-one-out, length $\ge 3$).}
\label{tab:flickr-ours}
\begin{tabular}{lccccc}
\toprule
Method & Toronto & Osaka & Glasgow & Edinburgh & Melbourne \\
\midrule
Random        & 0.298 & 0.301 & 0.301 & 0.270 & 0.227 \\
PoiPopularity & 0.443 & 0.413 & 0.510 & 0.439 & 0.320 \\
Markov (greedy) & \textbf{0.504} & \textbf{0.421} & \textbf{0.587} & 0.449 & 0.333 \\
MarkovPath (beam) & 0.528 & 0.398 & 0.543 & \textbf{0.452} & \textbf{0.346} \\
Pointer (learned, beam~3) & 0.431 & 0.399 & 0.489 & 0.414 & 0.312 \\
\bottomrule
\end{tabular}
\end{table}

\paragraph{Comparison to the literature.} Table~\ref{tab:flickr-pub} lists the directly
comparable published pairs-F1 (same protocol, same $0$--$1$ scale). Our classical line lands
between the published classical methods (PoiRank, Rank+Markov) on several cities, while the
neural SOTA (DeepTrip, CTLTR, SelfTrip, AR-Trip) reaches $0.66$--$0.85$ using machinery our
light model omits: self-supervised and contrastive pre-training, trajectory augmentation, and
adversarial training. There is no single SOTA across all cities: AR-Trip leads on
Toronto, Glasgow and Edinburgh, while SelfTrip is best on Osaka.

\begin{table}[ht]
\centering
\caption{Published, directly-comparable itinerary pairs-F1 on the Flickr cities (same Chen-2016
protocol). ``cl.''~=~classical, ``ne.''~=~neural.}
\label{tab:flickr-pub}
\begin{tabular}{llcccc}
\toprule
Method & Fam. & Toronto & Osaka & Glasgow & Edinburgh \\
\midrule
PoiRank~\cite{chen2016learning}      & cl. & 0.518 & 0.511 & 0.548 & 0.432 \\
Rank+Markov~\cite{chen2016learning}  & cl. & 0.512 & 0.486 & 0.545 & 0.444 \\
DeepTrip~\cite{gao2019deeptrip}      & ne. & 0.748 & 0.755 & 0.782 & 0.660 \\
CTLTR~\cite{zhou2021ctltr}           & ne. & 0.748 & 0.719 & 0.763 & 0.681 \\
SelfTrip~\cite{gao2022selftrip}      & ne. & 0.835 & \textbf{0.851} & 0.818 & 0.779 \\
AR-Trip~\cite{shu2024artrip}         & ne. & \textbf{0.839} & 0.828 & \textbf{0.820} & \textbf{0.808} \\
\bottomrule
\end{tabular}
\end{table}

\subsection{Personalization ablation}
\label{sec:sol-results-pers}

Because a per-user embedding is the system's only personalization mechanism, we isolate its
effect with a controlled ablation on the Flickr pointer (user embedding on vs.\ off; dimension
64, 30 epochs, leave-one-out). Table~\ref{tab:pers} shows the effect is \textbf{mixed}: it helps
clearly on Glasgow ($+0.038$ pairs-F1) but is roughly neutral-to-slightly-negative on Osaka and
Toronto. The honest reading is that, under leave-one-trajectory-out on these small datasets, most
users contribute too few trajectories for an identity embedding to be reliably estimated
(a cold-start effect); personalization helps only where user signal is both present and
consistent. The two larger cities (Edinburgh, Melbourne) are deferred to a GPU run.

\begin{table}[ht]
\centering
\caption{Personalization ablation: pointer pairs-F1 \emph{without} vs.\ \emph{with} the user
embedding (Flickr, length $\ge 3$, dim 64, 30 epochs).}
\label{tab:pers}
\begin{tabular}{lccc}
\toprule
City & no user emb. & + user emb. & $\Delta$ \\
\midrule
Glasgow & 0.451 & 0.489 & $+0.038$ \\
Osaka   & 0.442 & 0.435 & $-0.007$ \\
Toronto & 0.440 & 0.423 & $-0.017$ \\
\bottomrule
\end{tabular}
\end{table}

\subsection{Discussion}
\label{sec:sol-results-discussion}

The experiments support three honest conclusions, one per research question.
\textbf{(RQ1)} A learnable user embedding is the realized personalization mechanism; its benefit
is real but conditional (Glasgow $+0.038$), not universal---tourism's cold-start sparsity limits
identity-based personalization, motivating the persona/content-based direction as future work.
\textbf{(RQ2)} The spatial/temporal gap features $(\Delta d,\Delta t)$ and the hybrid POI graph
give the engine an honest LSTM/STGCN-tier accuracy ($\mathrm{HR@}1=0.187$); richer context such
as weather and time-of-day is a clear, but unrealized, extension.
\textbf{(RQ3)} A frozen next-POI engine can be decoded into coherent itineraries, and doing so
(Frozen Rollout) \emph{beats} a purpose-trained pointer because of its much denser supervision;
beam search alone cannot close the residual order gap (set-F1 $\gg$ pairs-F1), which points to
richer, constraint- and schedule-aware decoders (e.g.\ DLIR-style \cite{halder2025drlir}) as the
way forward. The Flickr Benchmark confirms that all of these conclusions sit on the literature's
own scale and that our harness reproduces Chen 2016, while honestly placing our light models
below the self-supervised neural SOTA.

% ============================================================
\section{Chapter Summary}
\label{sec:sol-summary}
% ============================================================

This chapter described the system as actually built: a context-aware next-POI engine (a GCN over
a hybrid POI graph, $(\Delta d,\Delta t)$ gap context, a user embedding, and a GRU scorer),
decoded into itineraries by Frozen Rollout and compared against an end-to-end Trained Pointer.
On Foursquare NYC the engine reaches an honest $\mathrm{HR@}1=0.187$, and---counter-intuitively
---decoding the frozen engine beats the integrated pointer ($0.289$ vs.\ $0.259$ pairs-F1).
On the Flickr Benchmark our methods land on the published $0.26$--$0.85$ scale with a
protocol-faithful harness (Random reproduces Chen 2016), our \texttt{Markov} baseline is strong,
and personalization is mixed. The next chapter draws these honest findings together and outlines
the unrealized extensions---weather/time context, persona-based cold-start, semantic curation,
and constraint-aware decoding---as future work.
